\chapter{Results, Analysis \& Discussions}

\section{Status as of \ReviewLabel}
At the time of this review, the project is in the architecture and
scaffolding stage: the system design (Chapter~3), data pipeline strategy,
and technology stack have been finalized, and the project's version-control
repository has been structured accordingly (data pipeline, machine
learning, backend, and frontend modules, each with defined interfaces).
Implementation of the live data pipeline, model training, and the
frontend application is in progress and had not yet produced evaluable
output at the time of this submission.

% TODO (Review 2 onward): replace this chapter's content with actual
% results as they become available -- e.g.:
%   - Sample of cleaned, backfilled historical data (row counts, date
%     coverage per state/commodity)
%   - Prophet vs. baseline MAE comparison table, per commodity/mandi
%   - Screenshots of the working dashboard / prediction page
%   - Any deviations from the Chapter 3 design and why

\section{Preliminary Data Analysis}
A sample single-day national Agmarknet export (9{,}584 records across 25
states, 271 districts, 642 markets, and 188 commodities) was analyzed to
validate assumptions made during architecture design. This confirmed, in
particular: (a) the field-name encoding artefact described in
Section~3.2.3 is present in bulk exports and not only the API preview;
(b) the variety/grade duplication pattern occurs in practice (e.g.\ two
distinct modal prices recorded for \textit{Arecanut} at the same market
on the same day); and (c) commodity reporting coverage varies
meaningfully by state -- for example, an originally assumed commodity
(Soyabean) had zero recorded rows for Maharashtra on the sampled day,
which directly motivated the shift to a data-driven, per-state top-crop
selection methodology (Section~3.1) rather than an assumed fixed crop
list.

\section{Discussion}
This preliminary analysis illustrates the value of validating
architectural assumptions against real data before implementation begins
-- several design decisions in Chapter~3 (state-based data organization,
data-driven crop selection, variety-deduplication logic) were refined
directly as a result of this analysis rather than decided purely from
documentation.
